\section*{第 \ref{ch:g4:magnets-and-compass} 章——磁铁与指南针}

\begin{solution}{exo:g4:magnets-and-compass:1}
\omterm{def:g1:magnets:magnet}{磁铁}的两个\omterm{def:g4:magnets-and-compass:poles}{磁极}是两端磁性最强的地方：\omterm{def:g4:magnets-and-compass:poles}{北极}（N，
常用红色标示）和\omterm{def:g4:magnets-and-compass:poles}{南极}（S）。回形针
聚集在两极——即两端——而不是中间。
\end{solution}

\begin{solution}{exo:g4:magnets-and-compass:2}
异名\omterm{def:g4:magnets-and-compass:poles}{磁极}相互\omterm{def:g1:magnets:magnet}{吸引}，同名\omterm{def:g4:magnets-and-compass:poles}{磁极}相互排斥。
发出“啪”的一声，说明两个异名\omterm{def:g4:magnets-and-compass:poles}{磁极}正相对，即 N 极朝向 S 极。
\end{solution}

\begin{solution}{exo:g4:magnets-and-compass:3}
N 极对 N 极：两者相互推开（同名\omterm{def:g4:magnets-and-compass:poles}{磁极}相斥）。N 极对 S 极：
两者相互拉近并吸合在一起，发出“啪”的一声。
\end{solution}

\begin{solution}{exo:g4:magnets-and-compass:4}
因为有一块巨大的\omterm{def:g1:magnets:magnet}{磁铁}在拉动它的\omterm{def:g4:magnets-and-compass:poles}{磁极}，使其
指向南北方向：地球本身就是一块\omterm{def:g1:magnets:magnet}{磁铁}，
\omterm{def:g4:magnets-and-compass:poles}{磁极}位于地球的南北两端附近。
\end{solution}

\begin{solution}{exo:g4:magnets-and-compass:5}
北、东、南、西。面朝北方时，东边在你的右手边。
\end{solution}

\begin{solution}{exo:g4:magnets-and-compass:6}
指针是一块小\omterm{def:g1:magnets:magnet}{磁铁}，它服从\emph{最近}的强大\omterm{def:g1:magnets:magnet}{吸引}。
附近的\omterm{def:g1:magnets:magnet}{磁铁}或大块铁器的\omterm{def:g1:magnets:magnet}{吸引}压过了远处地球这块\omterm{def:g1:magnets:magnet}{磁铁}，
使指针大致指向北方。
\end{solution}

\begin{solution}{exo:g4:magnets-and-compass:7}
用\omterm{def:g1:magnets:magnet}{磁铁}的一极沿同一方向摩擦指针 20 次，
每次摩擦之间把\omterm{def:g1:magnets:magnet}{磁铁}提起——指针便成为一块小\omterm{def:g1:magnets:magnet}{磁铁}。把它放在盛有
水的碗中\omterm{def:g1:floating-and-sinking:words}{漂浮}的软木片上：\omterm{def:g1:floating-and-sinking:words}{漂浮}的软木片相当于一个由水构成的\omterm{def:g3:levers-and-scales:lever}{支点}，
指针便会转动，
指向北–南方向。
\end{solution}

\begin{solution}{exo:g4:magnets-and-compass:8}
它指向北方。指针与正午\omterm{def:g1:light-and-shadows:shadow}{影子}一致的\omterm{def:g4:magnets-and-compass:compass}{指南针}
所指的方向是可信的——在晴天，
仅凭\omterm{def:g1:light-and-shadows:shadow}{影子}本身就能充当粗略的\omterm{def:g4:magnets-and-compass:compass}{指南针}。
\end{solution}

\begin{solution}{exo:g4:magnets-and-compass:9}
它在雾中、云下、无月的夜晚，以及远离任何海岸的地方都能工作
——指针对天空毫无要求。水手终于能在任何天气下，
在辽阔的大洋上保持航向。
\end{solution}

\begin{solution}{exo:g4:magnets-and-compass:10}
朋友皮带上的铁扣在近距离内压过了地球这块\omterm{def:g1:magnets:magnet}{磁铁}，
指针指向铁扣而非北方。她若信赖它，就会偏离方向——
错误的“北”位于她的左侧，她便会把路线向左拐，即偏向西方。
解决办法：先在离铁扣几步远的地方读取\omterm{def:g4:magnets-and-compass:compass}{指南针}，
然后再行走。
\end{solution}

\begin{solution}{exo:g4:magnets-and-compass:11}
这根短小的被摩擦过的指针有\emph{两}个\omterm{def:g4:magnets-and-compass:poles}{磁极}——它的两端
表现为 N 极和 S 极（让两根这样的指针\omterm{def:g1:floating-and-sinking:words}{漂浮}，它们会首尾相接，
异名端相互\omterm{def:g1:magnets:magnet}{吸引}）。因此缩短\omterm{def:g1:magnets:magnet}{磁铁}并不会去掉\omterm{def:g4:magnets-and-compass:poles}{磁极}，把\omterm{def:g1:magnets:magnet}{磁铁}对半切开
只会得到两块完整的短\omterm{def:g1:magnets:magnet}{磁铁}：靠切割寻找单独一个\omterm{def:g4:magnets-and-compass:poles}{磁极}的探索
永远不会结束——每切一次，
都会重新得到两个\omterm{def:g4:magnets-and-compass:poles}{磁极}。
\end{solution}
